\documentclass[11pt,a4paper]{article}
\usepackage[T1]{fontenc}
\usepackage[utf8]{inputenc}
\usepackage[brazil]{babel}
\usepackage{lmodern}
\usepackage[margin=2.3cm]{geometry}
\usepackage{booktabs}
\usepackage{amsmath,amssymb}
\usepackage[colorlinks=false,hidelinks]{hyperref}
\usepackage{enumitem}
\usepackage{fancyhdr}

\setlist{itemsep=2pt,topsep=3pt}
\pagestyle{fancy}\fancyhf{}
\fancyhead[L]{\small Perguntas prováveis --- UC day-ahead com VPL e VDC}
\fancyhead[R]{\small\thepage}
\renewcommand{\headrulewidth}{0.4pt}

% Pergunta destacada + resposta em parágrafo normal
\newcounter{perg}
\newcommand{\pq}[1]{\medskip\refstepcounter{perg}%
  \noindent\textbf{\theperg.\ #1}\par\nopagebreak\vspace{2pt}\noindent\ignorespaces}

\title{\vspace{-1.2cm}Perguntas prováveis do orientador e respostas\\
\large Co-otimização de Virtual Power Lines e Virtual Data Centers\\
no Unit Commitment day-ahead}
\author{Guilherme Steilein \\ \small Doutorado --- PPGEE-UFPR}
\date{\today}

\begin{document}
\maketitle
\thispagestyle{fancy}

\noindent
Trinta perguntas que considero prováveis, agrupadas por tema, com a resposta
curta que pretendo dar e o número que a sustenta. As que considero mais difíceis
são a 5, a 14, a 18 e a 26.

% =====================================================================
\section{Enquadramento e relação com a tese}

\pq{Por que juntar VPL e Virtual Data Center no mesmo modelo?}
Porque os dois aliviam o mesmo gargalo por caminhos opostos: a bateria move
energia no tempo, o data center move carga no espaço e no tempo. Cada um foi
estudado isoladamente, mas se eles se somam ou se anulam parcialmente é uma
pergunta em aberto, e a resposta muda o dimensionamento. É a pergunta que o
artigo responde.

\pq{O que exatamente é a métrica $I$ e por que essa definição?}
$I = \Delta_D - (\Delta_B + \Delta_C)$, isto é, o benefício do uso conjunto menos
a soma dos benefícios isolados. Negativo significa que juntos entregam menos que
a soma, ou seja, competem. É uma diferença de quatro objetivos de MILP, o que
exige cuidado numérico, mas fica na mesma unidade dos benefícios e pode ser
comparada diretamente com o ruído do gap do solver, o que uma razão não
permitiria.

\pq{A VPL do artigo é a VTL da tese?}
Não exatamente, e isso está explícito no texto. A tese define a VTL pela regra:
cada armazenamento opera permanentemente em carga ou descarga, conforme o estágio
de demanda líquida e o sentido do fluxo, e não presta outro serviço. No artigo
essa é a variante \emph{dedicada}. A formulação padrão é a \emph{emergente}, em
que as baterias são livres e o comportamento de linha virtual, se houver, emerge
da localização e dos preços.

\pq{Se o pareamento explica só 4\,\% do benefício, o conceito de VPL não perde sentido?}
Não, mas muda de lugar. Cada bateria sozinha vale 37,6 e 35,9\,k\$/dia, e o par
vale 76,6, então o pareamento acrescenta 4\,\%. A leitura é que quase todo o valor
vem de \emph{onde} o armazenamento está, e é justamente a localização que a tese
resolve. O que os números indicam é que, uma vez bem localizadas, as baterias
entregam o benefício sem coordenação explícita, ao menos num modelo determinístico
e sem contingências.

% =====================================================================
\section{Modelagem}

\pq{Por que a regra de chaveamento da tese \emph{custa} dinheiro em vez de agregar?}
Porque ela proíbe a bateria de arbitrar contra o sentido do fluxo nas horas
estagiadas. O benefício cai de 76,6 para 44,3\,k\$/dia na variante explícita e
para 40,0 na dedicada. Isso não é um defeito da regra: é o preço de um protocolo
executável pelo operador sem rodar um UC completo, com comportamento de
transferência garantido. Num modelo com contingências ou incerteza, essa garantia
pode valer mais que a arbitragem livre.

\pq{Sob a regra estrita, o que acontece com a bateria de montante?}
No dia de vento alto ela fica parada as 24 horas. O fluxo no corredor nunca
inverte, e com o sentido fixo a regra só permite que ela carregue nas horas de
alto uso e descarregue nas de baixo uso, o inverso do que a economia pede, porque
o excedente renovável de montante ocorre ao meio-dia. Nos outros três dias ela
opera de 5 a 12 horas.

\pq{Como o data center é modelado?}
Por tarefa, não por orçamento agregado de energia. Cada tarefa tem energia,
janela de execução, conjunto de sites viáveis e custo de migração assimétrico. No
PyPSA, cada tarefa é um barramento virtual alimentado por um \emph{link} de cada
site elegível, com eficiência $1/\mathrm{PUE}$, e terminado por um \emph{store}
cuja energia mínima salta no prazo. Isso separa por construção a flexibilidade
temporal da espacial.

\pq{Por que existe a restrição de rampa por tarefa, a Eq.~(10)?}
Ela veio do limite de rampa nativo do \emph{link} do PyPSA, mas não é redundante:
está ativa nas 168 instâncias com data center e em 8,7\,\% dos pares
tarefa--site por hora, um terço deles fora da borda da janela. Como altera a
região factível, entrou na formulação com parâmetro próprio, $\Delta^w_d$, e
justificativa física: o provisionamento é limitado por fluxo de migração, não só
no agregado.

\pq{As tarefas são divisíveis e interrompíveis. Isso não é otimista demais?}
É otimista, e está declarado como limite superior. Tarefas não interrompíveis
exigiriam binárias de ativação e dariam menos flexibilidade. Como o resultado
principal é que o data center \emph{compete} com a bateria, um modelo mais
restritivo tenderia a reduzir o benefício dele e, portanto, a competição, não a
inverter o sinal.

% =====================================================================
\section{Comparabilidade entre os casos}

\pq{Como garantir que os casos A/B e C/D são comparáveis?}
Nos casos sem flexibilidade os data centers continuam no sistema, como carga
plana, consumindo exatamente a mesma energia diária. Sem essa equalização, a
comparação mediria o valor de \emph{retirar carga} e não o valor da
flexibilidade. Esse cuidado vale cerca de 40\,\% do valor atribuído ao data
center: o primeiro código tinha as energias desiguais e inflava $\Delta_C$.

\pq{O requisito de reserva era o mesmo nos quatro casos?}
Não era, e isso foi testado. Em A e B a carga flexível entra como carga plana e
conta no requisito; em C e D ela passa por \emph{links} e não conta. Igualando o
requisito nos quatro casos, $\Delta_C$ cai de 24,3 para 23,9\,k\$/dia, dentro do
ruído, e $I$ vai de $-1{,}3$ para $-1{,}4$. A redução de partidas, porém, cai de
41 para 38\,\%: cerca de 3 pontos vinham da reserva mais folgada.

\pq{Por que o caso base não deixa os recursos flexíveis proverem reserva?}
Por conservadorismo. Com $\phi=1$ o efeito é de 0,01 a 0,2\,\%, praticamente
nulo, porque a reserva raramente é restritiva depois que o data center já achatou
a demanda líquida. Além disso, essa contribuição é um \emph{headroom proxy}: o
modelo não verifica se a tarefa interrompida ainda cumpre o prazo depois do
evento, nem se o estado de carga sustenta o compromisso da bateria.

\pq{E o baseline plano dos casos A e B, não favorece o data center?}
Favorece, e há a sensibilidade. Com as mesmas tarefas rodando no site nativo por
ordem de prazo, em vez de carga plana, $\Delta_C$ cai um terço, de 24,3 para
15,9\,k\$/dia. A interação muda de magnitude, de $-1{,}3$ para $-2{,}4$, mas não
muda de sinal em nenhum dos quatro dias. Mantive o plano como principal e este
como sensibilidade.

\pq{Por que o estado de carga inicial das baterias é livre?}
Porque o ciclo é fechado e o modelo escolhe o nível. É uma liberdade que o
operador real não tem. Ancorando em 50\,\%, o benefício da VPL cai 11\,\%, de
76,6 para 68,0\,k\$/dia, mas a interação fica \emph{mais} negativa, de $-1{,}3$
para $-1{,}7$. Ou seja, a hipótese é otimista quanto ao valor da VPL e
conservadora quanto à conclusão.

% =====================================================================
\section{Dados, sistema teste e método numérico}

\pq{De onde vem o sistema teste e a configuração estressada?}
RTS-GMLC, 73 barras em três áreas, com a modificação da tese e do artigo de 2026:
mais 40\,\% de carga na Área 2 e geração renovável equivalente na Área 3, nos nós
candidatos do planejamento. O corredor de estudo é o 318--223, de 500\,MW.

\pq{Por que quatro dias e não o ano inteiro?}
Por custo computacional e por escolha metodológica: são os quatro dias que a tese
usa, selecionados por critérios de demanda líquida. Mas eles são extremos, então
rodei também doze dias úteis medianos, um por mês.

\pq{E os dias mensais confirmam?}
Sim, com magnitude menor, como esperado. O custo cai 4,6\,\%, o curtailment
15\,\% e as partidas 47\,\%, contra 6,2, 12 e 41\,\% nos quatro dias. A interação
é negativa em 9 dos 12 dias, com média de $-2{,}7$\,k\$/dia. O sinal se mantém ao
longo do ano.

\pq{Por que 0,05\,\% de gap? Não é rigor excessivo?}
Não, é o que torna o resultado mensurável. $I$ é a diferença de quatro objetivos,
e os erros de otimalidade se somam em vez de se cancelar. Com a tolerância usual
de 0,5\,\%, cada caso custa cerca de 1,56\,M\$/dia, o erro admitido chega a
7,8\,k\$ por caso e a $\pm$15,6\,k\$ na soma dos quatro, maior que qualquer
interação medida. A 0,05\,\% o ruído cai para $\pm$1,6\,k\$.

\pq{O que teria acontecido se você tivesse ficado com 0,5\,\%?}
O sinal de $I$ seria inconclusivo em todas as combinações da grade, e chegou a
inverter entre as duas tolerâncias em cinco delas. A conclusão do trabalho não
existiria. É por isso que trato o ruído do gap como contribuição metodológica e
o reporto ao lado de cada interação.

\pq{Por que Gurobi e não HiGHS?}
Os dois aceitam 0,05\,\%; a diferença é o custo computacional. No caso A, o HiGHS
leva 80,5\,s contra 9,3\,s do Gurobi, 8,8 vezes mais. A licença acadêmica da UFPR
foi o que tornou viável rodar a grade inteira, 328 cenários, em algumas horas.
Mediana de 10\,s por caso; cinco instâncias pararam no limite de 1800\,s, com
gaps de 0,05 a 0,12\,\%.

% =====================================================================
\section{Resultados}

\pq{Qual é o mecanismo da competição, se os dois fazem coisas diferentes?}
Eles precisam do mesmo insumo escasso: a folga do corredor nas horas em que ele
não está saturado. A bateria de jusante só carrega com energia que atravessa o
corredor; a tarefa migrada para a Área 3 só tem valor porque consome ali a
energia que seria cortada. Quando um ocupa essa folga, sobra menos para o outro.
Nos dados, no caso conjunto o deslocamento temporal do data center e a capacidade
virtual da VPL ficam ambos menores que nos casos isolados.

\pq{Por que a competição é mais forte a 30\,\% de renovável e não a 70\,\%?}
Porque a 30\,\% o corredor satura em poucas horas, quatro em média, e os dois
disputam a mesma janela curta: o benefício conjunto fica 11,1\,\% abaixo da soma.
A 50 e 70\,\% ele satura de 9 a 11 horas, há mais espaço para atuarem em horas
diferentes, e a interação encolhe sem mudar de sinal. A hipótese inicial era a
oposta, e os dados não a sustentam.

\pq{No porte base a interação está dentro do ruído. Isso não enfraquece o artigo?}
Enfraqueceria se eu afirmasse mais do que os dados mostram, e é por isso que o
resumo diz explicitamente que no dimensionamento-base o efeito é da ordem da
incerteza numérica. O sinal mensurável aparece com VPL grande: a 400\,MW a
interação é $-6{,}7$\,k\$/dia, ou 28\,\% do benefício do data center, e a 650\,MW
chega a $-16{,}3$, cerca de 67\,\%. A mensagem prática é para dimensionamentos
grandes.

\pq{A conclusão depende de você ter deixado as baterias livres?}
Não. Repeti a comparação a 400\,MW nas três leituras. A interação é $-6{,}7$ na
livre, $-3{,}6$ na explícita e $-3{,}0$ na dedicada, e nas duas variantes com
regra ela é negativa nos quatro dias, contra dois dos quatro na livre. A
magnitude cai pela metade, porque uma bateria amarrada captura menos da
oportunidade e por isso também disputa menos. O sentido não muda.

\pq{Por que o data center reduz mais partidas que a bateria?}
Porque ele achata a demanda líquida na hora certa, concentrando as tarefas
adiáveis no pico renovável do meio-dia e deslocando 14\,\% da energia de TI
adiável para outros sites. A bateria tem energia limitada a 4\,h e atua sobre o
excedente renovável, não sobre a rampa da noite. Nas médias, o data center corta
7,8 partidas contra 5,5 da VPL de 200\,MW, evitando metade do curtailment e um
terço do custo que ela evita.

\pq{Qual é a consequência prática para planejamento?}
Dimensionar armazenamento de corredor e carga flexível de forma independente
superestima o benefício conjunto: de 1 a 11\,\% no porte base e até dois terços
do benefício da carga flexível quando a bateria é grande. A métrica $I$ sai de
quatro rodadas day-ahead e é um número que um estudo de dimensionamento pode
incorporar.

% =====================================================================
\section{Limitações e continuidade}

\pq{Quais são as limitações que você declara?}
Modelo determinístico, sem segurança N-1; quatro dias representativos mais doze
mensais, não o ano completo; tarefas divisíveis e interrompíveis, o que é limite
superior da flexibilidade; data centers sintéticos em tamanho, PUE e perfil; e
economia operacional, sem custo de investimento. O trabalho responde ``quanto
vale operar'', não ``vale a pena instalar''.

\pq{Os data centers são sintéticos. Isso não fragiliza o resultado?}
Fragiliza a magnitude, não o sinal. O valor do data center depende da fração
flexível e da janela: $\Delta_C$ vai de 11,3 a 37,2\,k\$/dia quando $f_d$ passa
de 0,2 para 0,8, e de 8,0 a 33,2 quando a janela é dividida ou dobrada. Por isso
reporto a faixa em vez de um número único, e por isso o próximo passo natural é
dados reais de carga de data center.

\pq{Por que o resultado é contraintuitivo e por que devo acreditar nele?}
Porque a intuição diz que serviços diferentes se somam. O que os dados mostram é
que serviço diferente não implica recurso complementar quando o insumo escasso é
o mesmo. Testei sete hipóteses de modelagem que poderiam produzir o resultado
artificialmente: precisão do solver, condições iniciais das térmicas, baseline do
data center, custo de migração, estado de carga inicial, carga e descarga
simultâneas e simetria do requisito de reserva. Nenhuma inverte o sinal.

\pq{Qual é o próximo passo, UC estocástico ou sistema brasileiro?}
Os dois são defensáveis e eu gostaria da sua opinião. O estocástico testa se a
competição se mantém sob incerteza, que é a crítica mais provável de um revisor,
e conecta com a linha de robustez da tese. O sistema brasileiro com dados reais
tem mais apelo para a conferência na América Latina e resolve a limitação dos
data centers sintéticos. Minha inclinação é o estocástico, por ser a extensão
metodológica mais forte.

\end{document}
